\documentclass[../main.tex]{subfiles}
\begin{document}
%==========================================TIÊU ĐỀ TẬP========================================
\begin{table}[h]
  \begin{adjustwidth}{-1cm}{}
    \begin{tabular}{>{\centering\arraybackslash}p{6cm}|>{\raggedright\arraybackslash}p{12.5cm}}
      \multirow{5}{6cm}
      % Cột bên trái: ÉPISODE và số tập
      {\\[-40pt]
        \flaregothic\fontsize{20pt}{20pt}\selectfont \centering
        \color{eptype}ÉPISODE\color{black}\\
        \burbank\fontsize{72pt}{72pt}\selectfont
        \color{epnum}129
        \\[18pt]
        % \flaregothic\fontsize{14pt}{14pt}\selectfont
        % \color{epattr}BIENVENUE, \uline{\color{Banpire} Mel} !
        % \flaregothic\fontsize{18pt}{18pt}\selectfont
        % \color{epattr}ÉPISODE PERDU
        \color{black}
      }
       &
      % Dòng thứ nhất: tên tập tiếng Nhật
      {\fotskip\fontsize{18pt}{18pt}\selectfont \ruby{BAR}{ばーる}でイキんな
      }                                                                                                                \\[6pt]
      % Dòng thứ hai: tên tập tiếng Anh
       & {\cafeteria\fontsize{18pt}{18pt}\selectfont Bar FUBAR}                                                                                                                    \\[4pt]
      % Dòng thứ ba: tên tập tiếng Việt
       & {\vnmsans\fontsize{18pt}{18pt}\selectfont Quán bar bất ổn}                                                                                                                \\[8pt]
      % Dòng thứ tư: Nhân vật xuất hiện
       & {\sen{\fontsize{14pt}{14pt}\selectfont Track used: Chess Fortress  (\ruby{割}{かっ}\ruby{拠}{きょ}\ruby{盤}{ばん}\ruby{上}{じょう}　\ruby{逐}{ちく}\ruby{鹿}{ろく}\ruby{群}{ぐん}\ruby{雄}{ゆう})}}
      \\[6pt]
      % Dòng cuối cùng: Thời gian diễn ra của tập (thường sẽ là ngày phát hành)
       & {\continuum\fontsize{16pt}{16pt}\selectfont Ngày phát hành: 31/10/2021}                                                                                                   \\[8pt]
       & (dựa theo bản dịch của \href{https://www.facebook.com/mamisubteam}{MamiSubs - Mami Fansub})                                                                               \\[18pt]
    \end{tabular}
  \end{adjustwidth}
\end{table}
%=========================================TRUYỆN CHÍNH========================================
\color{black}

Ngày cuối cùng của tháng Mười. Mùa thu đã trôi qua hai tháng.

Lẽ ra, đây phải là thời điểm mà cả văn phòng \hll\ ngập tràn trong không khí Halloween: những quả bí đỏ to đùng được chạm khắc khuôn mặt kỳ quái, ánh đèn lờ mờ hắt qua lớp rèm mỏng, những câu chuyện kinh dị được kể, và những bộ đồ hóa trang bắt đầu xuất hiện từ sáng sớm cho đến đêm muộn. Ấy vậy mà, vào ngày hôm ấy, không có gì diễn ra như vậy cả.

Không có ai hóa trang thành phù thủy (dù công ty của họ đã có một nàng phù thủy như thế), cũng không có ai rải kẹo ngoài hành lang. Không một chút u ám ma mị nào của ngày lễ ma quỷ. Thay vào đó, nó chỉ là một ngày hỗn loạn, kiểu hỗn loạn mà nếu có hay không cũng chẳng thay đổi mấy điều, tức là vẫn y như mọi ngày bình thường khác ở nơi này.

Và sự hỗn loạn đó không nằm ở văn phòng, mà lại đến từ\dots\ một quán rượu.

Ngay bên cạnh văn phòng, có một quán bar nhỏ mới mở được vài hôm, tên gọi nghe tưởng như quán ăn gia đình: ``Quán Cơm Nắm (hay おむすびばー trong tiếng Nhật).'' Tên thì dễ thương và thân thiện, nhưng nội dung bên trong thì chưa chắc.

Kuon và Tokue hôm nay rủ nhau sang đó làm vài ly giải khát, với lý do chẳng lấy gì làm chính đáng hơn ngoài cái gọi là ``đổi gió.'' Cả hai cậu bạn đang đứng bên ngoài, tỏ vẻ hứng khởi như những người đàn ông sau giờ làm.

``\vie{Ông nói cái quán này mới mở hả, Kuon?}''

``\vie{Cũng phải được một tuần rồi đấy,}'' cậu Tổng đáp. ``\vie{Tôi với một người nữa trong công ty tôi thành khách quen ở đây luôn rồi.}''

Cậu bạn Hiến máu cười nhếch mép: ``\vie{Chẳng biết thế nào, nhưng ông Kuon mà rủ thì đi thôi. Biết đâu lại là chỗ nhậu ra trò.}''

``\vie{Cũng có thể,}'' cậu nhún vai.

Cánh cửa gỗ được đẩy ra, kêu *KÉT* một tiếng nho nhỏ. Bên trong là một không gian mà Tokue phải dừng lại vài giây để tiêu hóa hết vì choáng ngợp: Sàn, trần, tường, quầy, bàn ghế, tất cả đều bằng gỗ, và mang một màu nâu ấm của thời gian. Từ các chao đèn thả xuống là ánh sáng vàng sậm, gợi lên cảm giác như đang ở một quán saloon giữa miền Tây nước Mỹ thời những năm 1930.

\img{10-8a}

Có người thì mong đợi nó sẽ giống những quán bar phục vụ whiskey, cũng có người mong có người đập đàn banjo. Tất nhiên ở quán này, điều gì cũng có thể thực hiện được, nếu có vị khách nào yêu cầu, ít nhất là theo câu khẩu hiệu của chủ quán này.

``\vie{Không gian này\dots\ quen phết nhỉ?}''

``\vie{Nếu ông từng xem mấy phim cao bồi kiểu Mẽo, kiểu mấy ông chĩa súng bắn nhau píu-píu, thì chắc nhận ra ngay đấy.}''

Hai cậu bạn tiếp cận tới chiếc bàn rượu lớn, nơi mà Kuon bắt gặp cô chủ quán. ``A, Okayu-chan!''

Cô nàng Miêu từ phía trong bước ra, tay cầm chiếc khăn lông lau nhẹ vào chiếc ly. Nghe thấy giọng cậu gọi, cô ngẩng đầu lên, nhoẻn miệng cười niềm nở. ``A, Kuon-senpai\~{} Hôm nay anh còn dẫn theo bạn nữa hả?''

Kuon chỉ về phía cậu bạn xứ Nghệ: ``Tại cậu ta tò mò thôi em. Ông ấy cũng muốn đổi gió một tí.''

Tokue đảo mắt nhìn quanh, vẫn không thôi trầm trồ với không gian quán: ``Xem chừng quán cũng ngon nghẻ nha. Ấm cúng đấy chứ.''

Chưa kịp để chủ quán hỏi, Kuon đã ngồi xuống ghế và vẫy tay gọi lớn như một khách hàng thân quen: ``Chủ quán ơi, cho anh hai cốc bia!''

Cậu bạn tròn mắt, cảm thấy có gì đó không đúng: ``Ê, ê, g-gượm đã, đây là quán rượu mà!''

Mặc cho bản mặt có phần khó hiểu của cậu bạn, nàng Miêu cười tươi đáp lại: ``Có liền\~{} Hai cốc bia lớn luôn nhé?''

Nói rồi, không cần chờ cử chỉ của cậu, cô hí hửng đi vào trong. Cậu Hiến máu nhìn mọi thứ diễn ra trơn tru, thốt ra đúng hai từ: ``\vie{\dots Vãi chưởng.}''

Cậu Tổng bật cười, vỗ vai trấn an: ``\vie{Nói thật tôi cũng từng sốc y như ông đấy. Gọi là quán rượu cho sang mồm thôi, chứ ông muốn gọi nước suối đóng chai cũng được. Đừng hỏi vì sao.}''

Tới đây, cậu bạn Hiếu máu chỉ biết cạn lời mà cười trừ tạm tin tưởng. ``Đành vậy.''

Chưa đến một phút sau, Okayu bước ra với hai cốc bia lạnh vừa được rót từ vại, bọt trắng xóa tràn lên miệng ly. Cô đẩy nhẹ hai cốc tới phía họ, cười nhẹ nhàng: ``Của hai người đây\~{} Chúc hai người ngon miệng\~{}''

``\textbf{Cho bọn anh xin!}'' cả hai cậu chàng cùng đồng thanh. Nói rồi, họ nâng ly cụng nhẹ một cái như những người đàn ông đã quen cái kiểu tụ tập quán bia sau giờ tan tầm.

Hai cậu con trai vừa cụng cốc và uống một hơi, thì cánh cửa quán lại mở lần nữa. Cả ba người cùng lúc liếc về phía lối vào. Người bước vào là một cô gái tóc xanh nước biển, đôi mắt sáng như sao băng, dáng đi dứt khoát như thể đã quá quen với nơi này.

Không cần nhìn kỹ, Kuon đã biết ngay đó là ai.

Người đó tiến thẳng tới quầy, ngồi ngay cạnh Kuon, khoanh tay lên bàn và dõng dạc tuyên bố: ``Chủ quán, như mọi khi nhé.''

\img{10-8b}

``Có liền\~{}''

Nàng Miêu quay gót lấy chiếc bình shaker (bình lắc) rồi xóc xóc như một bartender chính hiệu, dẫu món cô chuẩn bị lại chẳng phải là cocktail.

Tokue ngả người, nói khẽ với Kuon: ``\vie{\hl{Đó chẳng phải là Sui-chan sao?}}''

Kuon gật đầu xác nhận: ``\vie{\hl{Ừ. Một trong hai khách quen tôi nói đấy.}}''

``\vie{\hl{Cô ấy cũng uống rượu á?}}''

``\vie{\hl{Ờm\dots{} Không hẳn,}}'' cậu Tổng cố gắng lựa lời để giải thích cho cậu bạn nhằm không khiến đối phương quá bất ngờ. ``\vie{\hl{Quán này không chỉ phục vụ đồ uống\dots}}''

Cậu còn chưa nói xong, thì Okayu đã đổ ra cốc không phải rượu, mà là một cục cơm nắm rong biển, với topping là tôm chiên vàng ruộm.

Cậu chỉ vào cốc đựng cơm, hoàn thành câu nói: ``\vie{\hl{\dots mà còn phục vụ cả đồ ăn nữa. Đôi lúc còn theo cách không ai ngờ tới, như ông vừa thấy.}}''

Cậu bạn há hốc mồm, tiếp nhận thêm thông tin có phần kỳ quặc đó, tới mức miệng cậu chỉ buông ra một câu ``Quào'' khẽ khàng.

Suisei cầm lấy cơm nắm, ăn một miếng to, rồi ngồi trầm ngâm, đôi mắt nhìn xa xăm như thể đang truy tìm điều gì đó vừa mới thoảng qua trong ký ức. Rồi cô đột ngột mở mắt, hơi nghiêng đầu: ``Em mới thay loại gạo hả?''

Okayu đáp nhanh: ``Không. Như mọi khi mà.''

Kuon ở bên cạnh, liền nói góp: ``Dạo gần đây kinh tế khó khăn mà chủ quán vẫn giữ được chất lượng thế này là đáng nể đấy.''

``Gạo đó cũng không đắt lắm đâu mà.''

``Vậy à\dots'' cả nàng Ca sĩ và cậu Tổng cùng gật đầu.

Cảm thấy tình hình có vẻ bắt đầu chuyển biến, Kuon lập tức quay sang thì thầm với Tokue: ``\vie{\hl{Từ giờ cứ để yên cho họ nha.}}''

``\vie{\hl{Họ là ai cơ?}}'' Tokue ngơ ngác.

``\vie{\hl{Rồi ông sẽ thấy.}}''

Vừa dứt lời, Suisei bất ngờ quay sang phía một cô gái tóc vàng bạch kim đã ngồi vào bàn rượu từ khi nào. Đó là Nene.

\img{10-8c}

``Thế, em bảo em muốn nói chuyện sao?'' người tiền bối bắt đầu cất lời.

``Em ư!?'' cô nàng đàn em giật mình, tay tự chỉ vào chính mình.

Ánh mắt nàng Ca sĩ bỗng sáng rực như đèn sân khấu bật sáng tối đa, khiến Nene và Tokue bất giác toát mồ hôi hột.

``\vie{\hl{Tường là họ sẽ nghiêm túc chứ\dots}}'' cậu bạn Hiến máu thì thầm.

``\vie{\hl{Cứ kệ họ đi.}}''

Ở phía bên kia chiến tuyến, nàng ``Hải cẩu'' gãi má, cố nghĩ ra một chủ đề để bắt đầu trò chuyện: ``À, ờ\dots\ Em ghét ớt chuông lắm ấy\dots''

``Và cả cà chua nữa,'' Kuon chêm thêm.

``Ờ đúng! Cả cà chua luôn!''

Nghe vậy, Suisei càng tỏ vẻ hứng thú hơn, đôi mắt cô sáng rõ như sao đêm. Cậu bạn Hiến máu ngơ ngác không hiểu vì sao một cuộc hội thoại về rau củ lại trở thành chủ đề bàn luận nóng bỏng, cho đến khi cậu Tổng đỡ lời: ``\vie{\hl{Hai người bọn họ đều ghét rau củ. Đặc biệt là ớt chuông và cà chua.}}''

``\vie{\hl{À\dots{} Ra thế.}}'' Nói rồi, cả hai cậu bạn tiếp tục ngồi từ xa theo dõi cuộc hội thoại có phần ngây ngô.

Nàng Ca sĩ chỉnh lại tư thế, hắng giọng một tiếng rồi nghiêm túc nói: ``Nếu em nhìn quá lâu vào quả ớt chuông, thì chính nó cũng sẽ chằm chằm nhìn lại em\dots'' Cô cắn nốt miếng cơm nắm dở như một người từng trải, ánh mắt ánh lên một sự u uẩn không rõ lý do. Đối diện, Nene thì chớp chớp mắt, không biết phải phản ứng sao.

Tokue quay sang Kuon, toát mồ hôi hột: ``\vie{\hl{Tôi cảm thấy có gì đó\dots{} bất an nha\dots}}''

Kuon chỉ nhấp một ngụm bia, rồi bình thản đáp lại: ``\vie{\hl{Cuộc vui mới chỉ bắt đầu thôi, anh zai ạ.}}''

Đúng lúc ấy, cánh cửa quán lại mở.

Lần này, bước vào là một bóng hồng tóc xanh lam, áo choàng trắng nhè nhẹ bay trong gió. Đó là Lamy, xuất hiện với nụ cười dịu dàng: ``Quán vẫn mở chứ?''

\img{10-8d}

Suisei vẫy tay, chỉ vào chiếc ghế mà cậu bạn Hiến máu đã ngồi ban nãy: ``Chờ em nãy giờ đó. Ngồi xuống đi.''

``Ể?''

``Chị biết nỗi phiền muộn của em, Lamy à.''

Nàng Yêu tinh tuyết ngồi xuống, ngơ ngác nhìn đàn chị của mình: ``Phiền muộn? Ý chị là\dots?''

Chưa để cô nói thêm một câu gì, nàng Ca sĩ liền nói: ``Là rượu. Đúng không?''

``Không hẳn đâu chị\dots''

Câu trả lời của Lamy khiến cho Tokue và Kuon ở bàn bên đều không nhịn được cười. Cậu bạn Hiến máu quay đi, khẽ cười: ``\vie{\hl{Vãi lờ\dots{} Quê chưa kìa.}}''

Cậu Tổng cố che miệng, nhưng vẫn không giữ được tiếng cười: ``\vie{\hl{Cái này gọi là tự tin quá đà đó.}}''

Dù dự đoán sai toét, nàng ``Sao băng'' vẫn gượng gạo, tỏ vẻ đạo lý bằng giơ ngón tay trỏ lên như một nhà triết học: ``Chị có lời khuyên cho em này\dots\ Đã uống xe thì không\dots''

Bất chợt, cô dừng lại, cố gắng tìm ra cho từ đúng. ``Ơ\dots\ Cái gì sau đó đấy nhỉ?''

Ở phía bên hai người con trai, họ vẫn tiếp tục cố nén cười.

``\vie{\hl{Kiểu này là quên bài rồi đấy.}}''

``\vie{\hl{Câu đó mà em còn không nói đúng được thì thôi đấy.}}''

Suisei cố chống chế, tiếp tục nói bừa những từ mà gắn liền với\dots\ chủ đề mà cô cùng Nene nói trước đó: ``Ờm\dots\ Đã uống xe thì không\dots\ thì ăn\dots\ cái gì ấy nhỉ?''

Cuối cùng, đích thân người được ``tư vấn'' đành đỡ lời hộ cô nàng tội nghiệp đó: ``Em đoán ý chị là `Đã uống xe thì không lái rượu bia' ạ\dots''

Câu nói có phần ngây ngô và sai thứ tự đó đến tai hai cậu chàng. Họ vừa nhịn cười vừa lắc đầu ngán ngẩm: ``\vie{\hl{Về học lại môn giáo dục công dân đi mấy đứa ơi\dots}}''

Mặc kệ ánh mắt thương cảm (và cả tiếng phì cười) của ba người, nàng Ca sĩ vẫn cam chịu số phận của mình bằng cách\dots\ gọi thêm một nắm cơm nữa từ Okayu, như để nuốt trôi sự bẽ bàng.

Okayu vẫn như mọi lần, tươi tỉnh hỏi: ``Đến liền nha\~{} cơm nắm tôm phô mai hay mentaiko?''

\ct

Cùng lúc đó, cánh cửa quán lại bật mở lần nữa.

Polka và Botan bước vào như một cặp song tấu quen thuộc, đồng thanh cất tiếng chào mang thương hiệu chủ quán: ``\textbf{Mogu mogu\~{} Okayu\~{}!}''

\img{10-8e}

Kuon ngẩng đầu lên, chào hỏi một cách thân thiện: ``A, cả hội đông đủ nhỉ?'' Sở dĩ cậu nói thế vì bộ tứ NePoLaBo lần này đã tề tựu tại quán rượu gỗ cũ kỹ, một cảnh tượng hiếm thấy bên công ty.

Ngay lập tức, Suisei như bị kích hoạt bản năng ``thích thể hiện'', quay sang tạo dáng kiểu ``ta đây biết tuốt'', khiến cả hai cậu con trai đều rơi vào trạng thái ngượng chín mặt thay cho cô.

Cô nàng đắc ý, tay chỉ về phía hai người: ``Chị biết hai em muốn nghe gì rồi!''

Lập tức, Kuon giơ tay ngăn cản, giọng có pha chút hóm hỉnh: ``Thôi, anh nghĩ là đủ rồi đó, Sủi-chan ơi. Không cần phải tỏ vẻ nguy hiểm đâu.''

Nàng Ca sĩ nheo mắt, giọng trầm đặc lại, toát ra một thứ sát khí hăm dọa: ``Im mồm. Em xiên anh bây giờ đấy.''

Trái lại, cậu chi nhún vai, đáp lại đầy tỉnh bơ như đã từng trải qua chuyện này: ``Trúng tim đen nên là phát cáu hả? Bài đó xưa như Diễm rồi.''

Cô nàng hậm hực, rõ ràng không vừa lòng. Cô hắng giọng, rồi đọc một câu nghe như thơ trừu tượng: ``Ao cũ, con ếch nhảy vào, vang tiếng nước xao\dots''

Cô mỉm cười mãn nguyện như thể vừa truyền đạt một chân lý triết học\dots\ nhưng trước mặt cô là hai gương mặt ngơ ngác không kém gì khán giả xem chương trình văn nghệ thiếu nhi mà không có bối cảnh.

``Thôi anh xin. Câu cú chả rõ nghĩa gì mà còn đòi làm nhà triết gia nữa,'' Tokue buột miệng, mặt không giấu nổi vẻ giỡn cợt.

``Thật. Anh còn thấy Sư tử với Cáo sa mạc mặt đơ kìa,'' Kuon chỉ vào hai nét mặt chưng hửng, Botan thậm chí còn nhướng mày vì không hiểu đàn chị đang nói cái gì.

Thấy không ai hiểu được ``chiều sâu nghệ thuật'', nàng ``Sao băng'' gầm lên: ``\textbf{Thôi nào!!}''

\begin{figure}[H]
  \centering
  \begin{subfigure}[t]{0.45\textwidth}
    \centering
    \includegraphics[width=8cm]{../../../../Image/Book?/10-8f1.png}
    \caption{Polka và Botan ngơ ngác nhìn Suisei sau khi cô đọc xong bài thơ.}
  \end{subfigure}
  \quad
  \begin{subfigure}[t]{0.45\textwidth}
    \centering
    \includegraphics[width=8cm]{../../../../Image/Book?/10-8f2.png}
    \caption{Suisei nổi cáu vì không ai hiểu được những ``tầng lớp'' mà cô nói, vốn dĩ đã hoàn toàn vô nghĩa từ đầu.}
  \end{subfigure}
\end{figure}

Cô bất ngờ đứng bật dậy, tay tạo dáng kiểu biến hình, miệng lẩm bẩm như đọc lời thoại trong phim siêu nhân:

\begin{center}
  Ta sẽ nắm lấy mọi lo lắng, sử dụng chúng làm cơm nắm và ăn hết chúng!\\
  Duy ngã độc tôn.\\
  Đàn gảy tai trâu.\\
  Ta đây chính là Nữ Hoàng Cố Vấn!
\end{center}

Một chiếc vương miện rơi xuống từ hư không, gắn ba chữ to: Trợ Ngôn Vương (助言王) ở mặt trước.

\img{10-8g}

Kuon và Tokue cùng đổ mồ hôi hột trước sự xuất hiện của vương miện đó.

``\vie{Cái quán này có thả đạo cụ từ trần xuống luôn à\dots}'' Tokue ngước lên trần nhà, càng lúc càng khó hiểu hơn.

Kuon chỉ thở dài, chống hông lên bàn: ``\vie{Tôi chỉ mong là không phải do A-chan cài sẵn thôi\dots}''

Chưa kịp hoàn hồn, một giọng nói nữa vang lên: ``Có vẻ như không chỉ mình ngươi là nữ hoàng trong thị trấn này nhỉ.''

Cả quán ngước nhìn lên, điếng người trước sự xuất hiện của \emph{cô ấy}.

Polka và Botan thốt lên cùng lúc: ``Ủa!? Noel-senpai trên trần kia kìa!''

Thì ra là Noel, người đang đu bám trên xà gỗ bằng tay không. Cô đu mình xuống, vương miện trên đầu lấp lánh với dòng chữ: Não Cân Vương (脳筋王).

Suisei sững sờ khi lại có thêm một đối thủ tự xưng là ``vua'' khác: ``N-Nhà ngươi là\dots!''

Chỉ trừ hai cậu bạn vẫn còn đang hoang mang không biết chuyện gì đang diễn ra.

``Ờ\dots\ mọi người sao ngạc nhiên vậy?'' Kuon gãi đầu gãi tai, khó hiểu với cái tên ấy.

``Thế mà cũng là\dots\ vinh dự được à?''

``Chả trách người ta gọi em ấy là PON Knight.''

Ngay khi cả quán còn đang bối rối thì\dots\ *RẦM!* cửa quán bật mở lần nữa.

Kanata bước vào trong bộ dạng ngầu lòi, vừa đi vừa chống hông như đang chuẩn bị ``mổ bò đãi khách''.

Trên đầu cô là chiếc vương miện đính dòng chữ mạ vàng: PP Đồ Tể ({\fotskip{PP}王}).

``Người ta gọi ta là Nữ Vương Protein.''

Chưa kịp hết bàng hoàng, Luna ở kế bên Lamy cũng đứng dậy, tóc phấp phới như vừa có quạt thổi từ đâu đó. Cô mỉm cười, đội trên đầu một chiếc vương miện màu tím, đề: ``Una\~{} Vương (んな王).''

``Vua chuá chi mà chẳn có em dzô, kì dzạ-nanora\~{}?''

Và để hoàn tất cơn ác mộng cho hai cậu con trai, một giọng nói cuối cùng vang lên từ phía sau cây cột trong quán. ``Các ngươi có biết ai là vua thật sự không?''

Fubuki. Nàng Cáo trắng chậm rãi bước ra, để lộ chiếc vương miện lấp lánh kim tuyến: FUBUKING (フブキング).

\begin{figure}[H]
  \centering
  \begin{subfigure}[t]{0.45\textwidth}
    \centering
    \includegraphics[width=8cm]{../../../../Image/Book?/10-8h1.png}
    \caption{Vua Não Nho Noel.}
  \end{subfigure}
  \quad
  \begin{subfigure}[t]{0.45\textwidth}
    \centering
    \includegraphics[width=8cm]{../../../../Image/Book?/10-8h2.png}
    \caption{PP Đại Đế Kanata.}
  \end{subfigure}
  \quad
  \begin{subfigure}[t]{0.45\textwidth}
    \centering
    \includegraphics[width=8cm]{../../../../Image/Book?/10-8h3.png}
    \caption{Una\~{} Vương Luna.}
  \end{subfigure}
  \quad
  \begin{subfigure}[t]{0.45\textwidth}
    \centering
    \includegraphics[width=8cm]{../../../../Image/Book?/10-8h4.png}
    \caption{Fubuking Cáo trắng.}
  \end{subfigure}
  \caption{Lần lượt các thành viên khác đều xuất hiện với vương miện đề những chức vụ kỳ quặc.}
\end{figure}

Cậu Tổng toát mồ hôi hột nhìn mọi người: ``Ờm\dots\ thực sự trong này có bao nhiêu `vua' tự xưng thế?''

Cậu bạn Hiến máu cũng tỏ vẻ lạc lõng không kém: ``Sao tôi thấy giống thời kỳ loạn 12 sứ quân thời Đinh vậy?''

Nhưng sự bất ngờ vẫn chưa kết thúc ở đó.

Okayu, cô chủ quán rượu, người tưởng như đang đứng ngoài cuộc, đột nhiên xoay người lại, gật gù như đã dự tính từ trước. ``Có vẻ mọi người đã ở đây đầy đủ rồi nhỉ?''

Trên đầu của cô là một chiếc vương miện cùng loại với họ, đề rõ ba chữ: Tuyên Ngôn Vương (宣言王).

\img{10-8i}

Tất cả mọi người trong quán, kể cả những người không phải là ``vua'' đều há hốc mồm trước sự thật chấn động này.

Kuon há hốc mồm, tới mức làm rơi cả cốc bia xuống đất: ``Đến cả thông báo cũng là vua sao!?''

Còn Tokue thì chỉ có thể giơ hai tay ra hàng như tự nộp mình trước cảnh sát: ``Đến đây thì tôi chịu trận thật rồi đấy.''

Nàng Miêu nở nụ cười dịu dàng nhưng đầy quyền lực, dõng dạc tuyên bố với mọi người: Vậy thì ta, Tuyên Ngôn Vương, xin tuyên bố khai mạc\dots\ ``\textbf{Giải Đấu Nữ Hoàng \hll!}''

Toàn quán bỗng chốc vỡ oà.

Trên đầu Okayu, dòng chữ phát sáng: 「王位決定戦」(Giải Đấu Vương Quyền). Dưới chân các cô gái, mặt sàn gỗ tự động tách ra từng ô tròn như sân khấu của trò chơi đối kháng.

Khi mọi chuyện dần rẽ hướng hoàn toàn khác, cả hai cậu con trai chỉ có thể tần ngần nhìn các cô gái bắt đầu tranh giành nhau.

Họ chỉ có thể cùng nhau đồng thanh: ``\vie{Biết thế hôm nay ở lại làm đồ án trường cho xong\dots}''

%==========================================TRUYỆN PHỤ=========================================
\omk

Một rừng thì không thể có hai hổ, nhưng trong cái quán rượu toàn gỗ mang tên Cơm Nắm, dường như rừng lại mọc thêm MƯỜI con. Cùng lúc.

Các ``nữ hoàng'' tự xưng thi đấu đủ loại trò: từ ăn nắm cơm không dùng tay, đọc thơ bằng giọng theo chỉ định, trả lời trắc nghiệm văn hóa đại chúng, cho tới cả đấu súng nước kiểu cao bồi Mỹ.

Từng trò ngớ ngẩn đó, hết cuộc thi này tới trò chơi khác, mỗi cái đều trôi qua trong tiếng cười và tiếng đổ vỡ của một cái gì đó - có lẽ là vì thất vọng, hoặc là quá vui sướng, hoặc là quá cay cú vì thua.

Thậm chí, bốn thành viên của Thế hệ thứ Năm lại trở thành những MC bất đắc dĩ, đến mức họ đang nghĩ thầm: ``\textit{Tại sao chúng ta lại ở đây chứ?}''

Tới vòng thi toán học, Lamy đọc lớn một câu hỏi: ``42 nhân với 3 là bao nhiêu?''

Noel ngập ngừng một hồi lâu, không nghĩ ra được đáp án. Ngay tức khắc, Suisei bật dậy bấm chuông, rồi ráo hoảnh đáp: ``\textbf{126!}''

``Chính xác! Suisei-senpai giành thêm một điểm!''

Nàng Ca sĩ rạng rỡ giơ tay chiến thắng, miệng hô ``Tôi là toán vương!'' trong khi những người khác (và cả khán giả nam đang ăn nhậu ở bàn bên) đều nhăn mặt vì mức độ ``toán mẫu giáo'' của câu hỏi.

Bên phía khán giả, chiếc bàn to nhất quán được để nguyên làm khu vực dành riêng cho hội người ngoài cuộc,'' tức là Kuon, Tokue và đám bạn đại học của cậu, gồm Fukuyasu, Yamazu, Miharu và Shunyo.

Okayu dù đang làm MC chủ quán kiêm phục vụ, nwhng vẫn chăm sóc bàn ấy đầy đủ như một bồi bàn chính hiệu. ``Mọi người có muốn gọi gì không?''

Fukuyasu nói: ``Cho anh xin thêm một cốc đá nữa nhé.''

``Có liền\~{}''

Yamazu cũng tiện thể: ``T-thêm ít nước lẩu được không?''

``OK luôn\~{}'' Cô Miêu nhanh nhẹn chạy ra chạy vào. Cảnh tượng ấy khiến Miharu cười khì, vỗ vai Kuon: ``\vie{Anh thế mà cũng chất đấy. Lại có cả quán để anh em mình tụ họp!}''

Cậu Tổng nhún vai đáp: ``\vie{Bất đắc dĩ thôi. Với cả cũng lâu rồi anh em mình mới gặp đủ mặt thế này. Gần như vậy.}''

``\vie{Tiếc là Hotaro phải đi làm thêm, không thì cũng đông đủ rồi,}'' Shunyo xì xụp bát nước, rồi chợt quay sang hỏi mọi người: ``\vie{À mà, lát nữa ai trả tiền đây?}''

``\vie{Cam-pu-chia đi,}'' Miharu đề xuất, tay cầm chén rượu. ``\vie{Mỗi ông góp một ít là xong.}''

Kuon đảo mắt nhìn hội ``nữ vương'' đang tiếp tục đánh vật với câu hỏi toán học, rồi nhìn lại mâm lẩu của nhóm mình.

Một ý tưởng có phần điên rồ chợt nảy lên. Cậu hướng về mọi người, mắt lấp lánh tinh ranh: ``\vie{Tôi có trò này hay hơn.}''

Mọi người cùng quay sang, tò mò về cái trò mà cậu chuẩn bị đặt ra.

``\vie{Bây giờ anh em đặt cược xem ai thắng trong cái cuộc thi\dots\ `Nữ hoàng tự xưng' này.
  Ai đoán sai thì trả hết tiền cho cả bàn.}''

Tất cả mọi người trong mâm đều bày tỏ cảm xúc khác nhau: người thì ủng hộ hết mình, kẻ thì nghi ngại nhìn nhau.

``\vie{Anh zai, có ổn không đấy? Bọn mình gọi đồ cũng đâu có ít\dots}'' cậu bạn Hiến máu nghi ngờ.

``\vie{Lâu rồi tôi chưa thấy Kuon bạo dạn như vậy đó,}'' cậu Đạo diễn bình luận.

``\vie{Tôi chỉ sợ các ông không mang đủ tiền thôi,}'' cậu Tổng nhấp tiếp một ngụm bia, cười khà khà. ``\vie{Tôi thì vô tư!}''

``\vie{Tôi mang đủ tiền rồi,}'' Fukuyasu kiểm tra lại trong tài khoản. Các bạn học khác của cậu cũng vậy.

Cậu chàng ``sad boy'' lớn giọng tuyên bố: ``\vie{\textbf{Chơi thì chơi, sợ quái gì hả anh em!}}''

Tất cả mọi người hướng về nhóm cô gái vẫn đang hăng say thi đấu, rồi như một chủ sòng bạc nắm quyền kiểm soát, Kuon thốt lên: ``\vie{Đặt cửa đi nào anh em! Tôi là tôi vote (bình chọn) cho Luna-chan!}''

Ở quầy rượu chính, Okayu nghe thấy, ngoái lại mỉm cười nhẹ nhàng như một người chứng kiến nhiều trò đời hơn cô tưởng.

Ngày hôm đó, đám con trai có được một bữa nhậu hiếm có, vừa có bia, có lẩu, vừa được xem một giải đấu buồn cười của các nữ hoàng tự xưng ``thi tài năng'' vì những lý do chẳng ai rõ.

Mỗi người mạnh ở một mảng (nếu nói cho lịch sự), còn nếu nói thẳng ra thì ai cũng yếu ở phần còn lại. Sau mười mấy vòng, họ đều bằng điểm nhau, không ai hơn ai.

Và như thể có cùng một suy nghĩ, cả nhóm đồng thanh: ``\textbf{Tất cả mọi người đều cùng hòa!}''

Nói theo cách nói giảm nói tránh là chẳng ai thắng, và chẳng ai thua.

Hội anh em cũng dựa theo kết quả đó mà chia tiền trả bữa như thường lệ, nhưng bất ngờ thay  Okayu lại mang ra một tập phiếu giảm giá. ``Vì Kuon-senpai là khách quen nên mọi người dùng mấy phiếu này đi nha\~{}''

Món quà bất ngờ từ cô chủ quán khiến cho nhóm những người đàn ông vừa mừng rỡ, vừa cảm kích cho cô lẫn cậu Tổng.

``Thật á? Cảm ơn nha Okayu-chan!'' Tokue mừng ra mặt. ``Lại tiết kiệm được ít tiền về quê rồi!''

``Được giảm giá luôn thì sao phải ngại!'' Fukuyasu hóm hỉnh nói, cười tít mắt.

Chuyện tưởng như đã kết thúc, thì hóa ra vẫn còn một điều ngạc nhiên khác dành riêng cho vị khách quen ấy.

Suisei đứng bật dậy, chỉ tay về phía bàn của hội Kuon, những người đang ăn nốt nồi lẩu. ``Mọi người! Giải đấu không có vua thực sự thì không trọn vẹn. Vậy thì\dots\ hãy bầu chọn Vua của \textit{hololive} hôm nay!''

Mọi người xung quanh ngơ ngác với ý tưởng đột ngột này, nhưng rồi từng người, từng người đồng thanh: ``Kuon-senpai!'' ``Kuoncha-senpai-nora!'' ``Tôi cũng bầu cho anh ấy luôn!'' ``Không có ai khác xứng đáng hơn!''

Thậm chí đến cô chủ quán cũng đồng tình với các ``nữ vương'' tự xưng ấy: ``Anh ấy trả tiền đều đặn và không bao giờ quên bo đó\~{}''

Những lời nói dồn dập chỉ đích danh cậu khiến cho bản thân Kuon giật mình, suýt làm rơi cả bát mì trên tay: ``Hả!? Ủa, ủa, mọi người\dots\ Không lẽ\dots''

Tất cả mọi người trong quán trong khoảnh khắc đó không một ai phản đối.

Cuối cùng, Hikawa Kuon, một cậu quản lý tận tâm với công việc và tốt bụng với mọi người, một cậu chàng chỉ đi uống rượu cho vui, bất đắc dĩ được phong là Vua \hll\ trong tiếng vỗ tay nổ trời và những tràng cười vỡ bụng của bạn bè.

``Ờm\dots\ cảm ơn mọi người? Anh-anh cũng không biết phải bày tỏ cảm xúc thế nào luôn\dots'' nhà vua nói, cười gượng trong khi khuôn mặt đỏ hoe.

Yamazu vỗ vai, cười lón: ``\vie{V-Vua bất đắc dĩ là có t-thật!}''

Shunyo thì hô vang: ``\vie{Mọi người cùng nổ trào pháo tay thật lớn cho nhà Vua mới nào!}''

Fukuyasu bụm miệng cười: ``\vie{Vua \hll\ nghe cũng oách đấy chứ!}''

Và thế là, trong một ngày tưởng như chỉ là cuối tháng Mười không có Halloween, một ông Vua mới được sinh ra, không phải bằng kiếm hay phép thuật, mà bằng\dots\ cơm nắm, rượu, và những tiếng cười.

Có thể đó chẳng phải là ngày đáng nhớ gì để được vinh danh vào sách lịch sử hay sách kỷ lục thế giới, nhưng với tất cả những người có mặt trong quán hôm đó, đó lại là một kỷ niệm vừa hài hước vừa đáng nhớ khác được mọi người lưu giữ mãi trong tâm mình.

\end{document}